\section{Puente 2: traer experiencias al Contexto mediante Episodic Retrieval}

La segunda ruta de puente no anticipa todas las conclusiones por adelantado, sino que preserva experiencias más completas. Cuando surgen nuevas preguntas, el sistema recupera episodios relevantes y los reintegra en el contexto para que el modelo utilice sus habilidades ICL existentes e interprete según el objetivo actual. Este método transforma "uso futuro desconocido" en "si se pueden encontrar las experiencias correctas durante la prueba".

\subsection{Memoria episódica de oracle: demostrar valor primero, reconocer luego que la recuperación no está resuelta}

El recuperador del curso es una memoria episódica de oracle: garantiza recuperar al menos una experiencia relevante, aunque mezcle varias distracciones. Este diseño divide intencionalmente el problema: primero se pregunta "si una experiencia relevante regresa al contexto, ¿se recupera la generalización latente?", en lugar de resolver simultáneamente embebido, índice, clasificación por relevancia y recencia junto con los controles de permisos.

\lecturefigure{slide-033.jpg}{Enfoque 2: recuperación episódica de oracle restaura experiencias relevantes durante la prueba}{Diapositiva oficial de Stanford CS25 V6 Lecture 06 p. 33}

Una vez definida el conjunto de recuperación $M_q$, las métricas comunes son
$$
\operatorname{Recall}@k=\frac{|M_q^{(k)}\cap R_q|}{|R_q|},
\qquad
\operatorname{Precision}@k=\frac{|M_q^{(k)}\cap R_q|}{k}.
$$
donde:
\begin{itemize}
  \item $R_q$ es el conjunto de experiencias relevantes para la pregunta $q$;
  \item $M_q^{(k)}$ son las primeras $k$ experiencias devueltas por el recuperador;
  \item la configuración de oracle garantiza un recall alto, pero permite que la precisión no sea perfecta;
  \item por ello, los experimentos prueban si el modelo puede utilizar al menos una memoria correcta entre las distracciones.
\end{itemize}

\begin{lstlisting}[language=Python,caption={Separación de responsabilidades en el experimento de recuperación episódica de oracle}]
memories = store_training_experiences(dataset)

def answer(query):
    relevant = oracle_pick_at_least_one(memories, query)
    distractors = sample_irrelevant(memories, count=4)
    context = shuffle(relevant + distractors)
    return frozen_model.generate(context=context, query=query)
\end{lstlisting}

\teachervoice{00:29:09--00:29:31, el docente llama directamente a este sistema de recuperación como "tramposa": tiene recall perfecto, aunque la precisión no lo sea. Este reconocimiento no debilita el experimento, sino que aclara que solo prueba una proposición condicional: si un episodio relevante regresa al contexto, efectivamente tiene valor.}

\subsection{Resultados: reversal y codebook se recuperan}

El foco de la comparación en las gráficas no es la condición forward, porque tanto el aprendizaje de parámetros como la recuperación episódica pueden manejar direcciones explícitas; lo verdaderamente diferenciador son la condición reversal y la condición de codebook latente. Un sistema con recuperación puede acercarse a una generalización flexible estilo ICL incluso al ver memorias distractores simultáneamente.

\lecturefigure{slide-034.jpg}{Beneficios de la recuperación episódica en las pruebas de reversal y codebook latente}{Diapositiva oficial de Stanford CS25 V6 Lecture 06 p. 34}

\begin{knowledgebox}{Lectura de la figura: condiciones explícitas para calibrar, condiciones latentes para distinguir mecanismos}
En ambos grupos de gráficas, primero se deben observar las condiciones explicit / trained para confirmar que el aprendizaje de parámetros y la recuperación pueden completar al menos las tareas básicas; luego se revisan las condiciones reversal o codebook latente. Si solo estas últimas marcan diferencia significativa, la evidencia respalda un uso flexible, no que la recuperación simplemente provee más datos de entrenamiento o un modelo base más fuerte.
\end{knowledgebox}

\lecturefigure{slide-035.jpg}{Conclusión: traer experiencias al contexto puede recuperar su valor latente}{Diapositiva oficial de Stanford CS25 V6 Lecture 06 p. 35}

\begin{importantbox}{La recuperación cambia el entorno computacional}
Los parámetros podrían retener hechos, pero no exponer de manera estable la dirección necesaria para completar la transformación actual; la recuperación episódica presenta las experiencias originales como tokens nuevamente, permitiendo que el modelo invoque los circuitos de razonamiento por contexto ya entrenados. No se escriben más hechos permanentemente en los parámetros, sino que se reconstruye un mundo local computable para el objetivo actual.
\end{importantbox}

\subsection{Relación con RAG: no solo buscar respuestas, sino también experiencias reinterpretables}

El RAG ordinario suele describirse como la recuperación de documentos que contienen respuestas; aquí nos interesa más recuperar experiencias \textbf{de las cuales se pueden deducir respuestas}. Por ejemplo, si un documento solo dice "X es padre de Y", pero la pregunta es "¿Quién fue padre de Y?", el recuperador no necesita almacenar una respuesta inversa pregenerada, sino que basta con devolver la relación original al contexto.

\begin{knowledgebox}{Diferencias entre recuperación episódica y RAG factual convencional}
\begin{tabularx}{\textwidth}{L{0.24\textwidth}X X}
\toprule
Dimensión & Objetivo común en RAG factual & Recuperación episódica en esta clase \\
\midrule
Objeto de recuperación & Párrafos con respuesta directa o similares & Experiencias originales relevantes al objetivo actual \\
Responsabilidad de generación & Extraer, sintetizar y citar hechos & Reinterpretar direcciones, combinar relaciones y cambiar objetivos \\
Riesgos clave & Documentos faltantes, orden incorrecto, citas falsas & Conservar suficiente detalle relacional además de la relevancia \\
Estado experimental & Evaluación con recuperadores reales & El curso aísla el valor de la recuperación usando oracle \\
\bottomrule
\end{tabularx}
\end{knowledgebox}

\teachervoice{00:30:41--00:31:24, el docente vincula esta tarea con la generación aumentada por recuperación, pero señala que el enfoque es permitir que el contexto realice razonamiento nuevamente sobre relaciones latentes difíciles de invocar en los parámetros, no solo consultar una respuesta factual.}

\begin{warningbox}{Los resultados de oracle no pueden convertirse directamente en promesas de producto}
Un sistema real debe resolver reformulación de consultas, deriva del embebido, contaminación por memoria a largo plazo, permisos y privacidad, validez temporal, experiencias conflictivas y presupuestos de recall/latencia. Si el recall disminuye, la arquitectura episódica conceptualmente correcta también fallará porque "la memoria no regresó".
\end{warningbox}

\subsection{Resumen de la sección}

La recuperación episódica evita enumerar offline todos los usos futuros posibles, devolviendo experiencias relevantes al contexto solo cuando surge un objetivo. El experimento con oracle demostró que esta interfaz computacional puede recuperar el valor latente del reversal y del codebook, aunque no resuelve un recuperador universal. La tercera ruta indaga aún más: si el modelo no accede a una memoria externa, ¿puede generar por sí mismo la información necesaria durante la prueba?


